\section{Sampler Audit: Nominal versus Realized Retention}
\label{sec:sampler}

Before presenting experiments, we audit what DSpar-style samplers actually retain, because the retention parameter $\alpha$ has been used in the literature---and in an earlier version of this work---with three mutually inconsistent meanings.

\paragraph{With-replacement sampling collapses retention.}
The originally published DSpar procedure~\cite{liu2023dspar} draws $q = \lceil \alpha m \rceil$ edge samples \emph{with replacement} from the distribution proportional to $s(e)$, reweights, and keeps the union of sampled edges.
Because high-score edges are drawn repeatedly, the number of \emph{distinct} retained edges is far below $\alpha m$: across our 17-network suite, nominal $\alpha = 0.8$ realizes only $0.33$--$0.52$ of the edges (Table~\ref{tab:sampler_audit}).
Results reported at ``$\alpha = 0.8$'' under this sampler are, in fact, roughly half-density experiments.

\paragraph{Probability clipping causes saturation.}
A natural no-replacement variant retains edge $e$ independently with probability proportional to its score, clipped at one.
Because a substantial fraction of low-degree edges clips (13--35\% in our measurements), the expected retention falls strictly below $\alpha m$ at every setting and \emph{saturates}: on \texttt{email-Eu-core} this sampler cannot retain more than $66\%$ of edges even at $\alpha = 1.0$.
True mild sparsification (e.g., retaining $90\%$ of edges) is unreachable under either published variant.

\paragraph{Calibration closes the gap.}
Definition~\ref{def:calibrated} instead solves for the proportionality constant $\lambda_\alpha$ by bisection so that expected retention equals $\alpha m$ exactly.
Measured realized retention under the calibrated sampler matches the nominal value to within $\pm 0.001$ at all tested settings ($\alpha \in \{0.7, 0.8, 0.9, 0.95\}$).
We verified our implementations of the with-replacement variant to be bit-identical to the original reference code before use.

\begin{table}[t]
\centering
\setlength{\tabcolsep}{4pt}
\caption{Sampler audit: nominal $\alpha$ versus realized edge retention $m'/m$, over all runs on
disk. The three samplers used in (or implied by) the pipeline are not interchangeable.
\texttt{paper} samples $\lceil \alpha m\rceil$ edges \emph{with} replacement, so distinct-edge
retention collapses to roughly $\alpha/2$ and never approaches $\alpha$ --- even at $\alpha=1.0$.
\texttt{probabilistic\_no\_replace} sets $p_e=\mathrm{clip}(s_e/\sum s\cdot\lceil\alpha m\rceil,0,1)$;
$13$--$35\%$ of edges hit the clip, so $\mathbb{E}[\text{kept}]<\alpha m$ always and the sampler
saturates below $0.7$. The calibrated sampler solves $\sum_e \min(1,\lambda s_e)=\alpha m$ by
bisection, giving $\mathbb{E}[\text{retention}]=\alpha$ exactly; it is the only variant for which the
nominal knob means what the formal definition says. ``\#'' is the number of network$\times$sampler
runs pooled into the range. Ranges are min--max of the per-network mean realized retention.
Sources: \texttt{PAPER\_RESTRUCTURE/audit/AUDIT\_FINDINGS.md},
\texttt{exp\_C/results\_summary.csv}, \texttt{exp\_B/results.csv}, \texttt{exp\_E/outcomes.csv},
\texttt{exp\_F/results.csv},
\texttt{PAPER\_EXPERIMENTS/results/exp4\_comprehensive/comprehensive\_alpha\_results.csv}.
Exp\_C additionally ran $\alpha\in\{0.7,0.95\}$ (\texttt{no\_replace}: $0.366$--$0.572$ and
$0.447$--$0.678$; calibrated: $0.699$--$0.700$ and $0.950$--$0.951$). The calibrated
$\alpha=1.0$ cell was never run, as it is the no-sparsification sentinel.}
\label{tab:sampler_audit}
\footnotesize
\begin{tabular}{llrrl}
\toprule
Sampler & nominal $\alpha$ & realized $m'/m$ & \# & source \\
\midrule
\multirow{3}{*}{\texttt{paper} (w/ repl.)} & 0.8 & 0.317--0.549 & 54 & exp4, exp\_B (both arms), exp\_F \\
 & 0.9 & 0.334--0.554 & 15 & exp4 \\
 & 1.0 & 0.350--0.589 & 15 & exp4 \\
\addlinespace[2pt]
\multirow{3}{*}{\texttt{prob.\_no\_replace}} & 0.8 & 0.400--0.619 & 6 & exp\_C \\
 & 0.9 & 0.432--0.660 & 6 & exp\_C \\
 & 1.0 & 0.662 & 1 & audit (email-Eu-core) \\
\addlinespace[2pt]
\multirow{3}{*}{calibrated ($\lambda$-bisection)} & 0.8 & 0.799--0.800 & 6 & exp\_C \\
 & 0.9 & 0.887--0.901 & 26 & exp\_C, exp\_E, exp\_F \\
 & 1.0 & -- & -- & not run \\
\bottomrule
\end{tabular}
\end{table}


\paragraph{Consequences.}
Three practices follow, which we adopt throughout.
First, realized retention $m'/m$ is reported alongside nominal $\alpha$ in every experiment.
Second, comparisons across sampler variants are made at matched \emph{realized} retention, not matched nominal $\alpha$.
Third, we reserve $\alpha = 1.0$ strictly as the no-sparsification baseline; note that under the with-replacement variant even $\alpha = 1.0$ removes roughly half the edges, so pipelines that treat it as a baseline must bypass the sampler explicitly.
